\section{Yöntem}
\label{sec:method}

Bu bölümde \fastgrnn{} hücresini, hücreye uygulanan üç aşamalı
sıkıştırma boru hattını ve modelin çarpıcısız bir MCU üzerinde
çalışabilmesini sağlayan LUT tabanlı aktivasyon yöntemini tanımlıyoruz.
Şekil~\ref{fig:pipeline}, tüm akışı özetler.

\subsection{\fastgrnn{} Hücre Formülasyonu}
\label{sec:method:cell}

$t$ anındaki giriş $\mathbf{x}_t \in \mathbb{R}^{d}$ ve önceki gizli
durum $\mathbf{h}_{t-1} \in \mathbb{R}^{H}$ verildiğinde,
\fastgrnn{} hücresi~\cite{kusupati2018fastgrnn} tek bir kapı
$\mathbf{z}_t$, aday güncelleme $\tilde{\mathbf{h}}_t$ ve iki
skalerli ara değerleme $\mathbf{h}_t$ hesaplar:
\begin{align}
\mathbf{z}_t &= \sigma(\mathbf{W} \mathbf{x}_t + \mathbf{U} \mathbf{h}_{t-1} + \mathbf{b}_z) \label{eq:z}\\
\tilde{\mathbf{h}}_t &= \tanh(\mathbf{W} \mathbf{x}_t + \mathbf{U} \mathbf{h}_{t-1} + \mathbf{b}_h) \label{eq:htilde}\\
\mathbf{h}_t &= \bigl(\zeta(\mathbf{1} - \mathbf{z}_t) + \nu\bigr)\, \tilde{\mathbf{h}}_t + \mathbf{z}_t \odot \mathbf{h}_{t-1} \label{eq:hupdate}
\end{align}
Burada $\odot$ eleman bazında çarpmayı, $\zeta, \nu \in (0,1)$ ise
sızıntılı-integratör davranışı ile standart kapılı güncelleme arasında
denge kuran \emph{öğrenilmiş skalerleri} gösterir. \fastgrnn{}'yu
ayıran nokta bu iki skalerli kapıdır: GRU ya da LSTM'deki üç veya
dört ağırlık çifti yerine, kapı ve aday güncelleme aynı ağırlık
çiftini $(\mathbf{W}, \mathbf{U})$ paylaşır. Böylece LSTM'ye yakın
ifade gücü çok daha az parametreyle elde edilir. HAR kurulumumuzda
$d{=}3$ (üç eksenli ivme),
$H{=}16$ ve giriş dizisi uzunluğu $T{=}128$ örnektir (50~Hz'de
2.56~s'lik bir pencere). $H{=}16$ seçimi Tablo~\ref{tab:hsize}'da
deneysel olarak gerekçelendirilmiştir: gizli boyutu $H{=}32$'ye
çıkarmak parametre sayısını yaklaşık üç kat artırdığı hâlde daha
düşük test F1 üretir ve eğitim uzadıkça sonuç daha da kötüleşir.
Bu davranış, daha küçük hücrenin veri kümesindeki yapıyı zaten
yakaladığı ve büyük modelin aşırı uyuma gittiği anlamına gelir.

\begin{table}[h]
\centering
\caption{\hapt{} üzerinde gizli boyut seçimi; tam ranklı
\fastgrnn{} (düşük rank yok, seyrekleştirme yok). $H{=}32$,
$H{=}16$'ya göre yaklaşık üç kat parametreye sahip olmasına rağmen
eşleşen her yapılandırmada daha kötüdür ve epoch~30'dan~50'ye
\emph{daha da} bozulur; bu durum aşırı uyuma işaret eder.
$H{=}16$ epoch~120'ye kadar iyileşmeye devam eder
(doğrulama-seçimli tavan F1~$=$~0.912; bkz.
Bölüm~\ref{sec:results:lsq}). $H{=}32$ 120~epoch'a
uzatılmamıştır; 30~$\rightarrow$~50 aralığındaki monoton bozulma,
daha uzun eğitimin farkı kapatma olasılığını son derece düşürür.
Bu nedenle sonraki tüm deneylerde $H{=}16$ kullanılmıştır.}
\label{tab:hsize}
\begin{tabular}{rrrrr}
\toprule
$H$ & Epoch & F1 & Doğruluk & Parametre \\
\midrule
\textbf{16} & 30  & 0.847          & 0.854          & \textbf{440}   \\
\textbf{16} & 50  & 0.861          & 0.869          & \textbf{440}   \\
\textbf{16} & 120 & \textbf{0.912} & \textbf{0.913} & \textbf{440}   \\
32          & 30  & 0.833          & 0.845          & 1{,}384         \\
32          & 50  & 0.830          & 0.834          & 1{,}384         \\
32          & 120 & ---            & ---            & 1{,}384         \\
\bottomrule
\end{tabular}
\end{table}

Aynı ağırlık çifti $(\mathbf{W}, \mathbf{U})$ hem~\eqref{eq:z}
hem de~\eqref{eq:htilde} içinde yeniden kullanıldığından, hücrenin
kısıtsız parametre sayısı şöyledir:
\begin{align}
\#\text{params}
 &= \underbrace{Hd}_{\mathbf{W}} +
    \underbrace{H^2}_{\mathbf{U}} +
    \underbrace{2H}_{\mathbf{b}_z,\mathbf{b}_h} +
    \underbrace{2}_{\zeta,\nu} \nonumber\\
 &= 48 + 256 + 32 + 2 = 338.
\end{align}
Aynı gizli boyuttaki standart bir LSTM, bu tür dört ağırlık matrisi
gerektirir ($\approx$\,1{,}280 parametre). Dolayısıyla
\fastgrnn{}, mimari düzeyde $\sim$$4\times$ üstünlükle başlar.
Aşağıdaki düşük rank ve seyrekleştirme aşamaları bunu bir büyüklük
mertebesi daha azaltır.

\subsection{Düşük Ranklı Ağırlık Ayrıştırması}
\label{sec:method:lowrank}

Giriş matrisi $\mathbf{W} \in \mathbb{R}^{H \times d}$ ve yinelemeli
matris $\mathbf{U} \in \mathbb{R}^{H \times H}$, iki ince matrisin
çarpımı olarak ayrıştırılır:
\begin{align}
\mathbf{W} &= \mathbf{W}_1 \mathbf{W}_2^\top, \quad \mathbf{W}_1 \in \mathbb{R}^{H \times r_w}, \mathbf{W}_2 \in \mathbb{R}^{d \times r_w}, \\
\mathbf{U} &= \mathbf{U}_1 \mathbf{U}_2^\top, \quad \mathbf{U}_1 \in \mathbb{R}^{H \times r_u}, \mathbf{U}_2 \in \mathbb{R}^{H \times r_u}.
\end{align}
$\mathbf{W}_1, \mathbf{W}_2, \mathbf{U}_1, \mathbf{U}_2$
çarpanları birlikte öğrenilir. Yinelemeli matris-vektör çarpımı
$\mathbf{U}\mathbf{h} = \mathbf{U}_1(\mathbf{U}_2^\top \mathbf{h})$
olarak değerlendirilir; böylece maliyet $H^2$ yerine $2 H r_u$
çarp-topla işlemine iner. $r_u \in \{4, 6, 8, 12\}$ değerlerini her
biri beş rastgele tohumla taradıktan sonra
(Bölüm~\ref{sec:results:lowrank}), $r_w = 2, r_u = 8$ seçilmiştir.

\subsection{İteratif Sert Eşikleme}
\label{sec:method:iht}

Dört çarpan matrisinin tamamını iteratif sert
eşikleme~\cite{blumensath2009iht} ile seyrekleştiriyoruz. Her eğitim
adımında, her ağırlık tensöründeki mutlak değerce en büyük $k$ girdi
korunur; diğerleri sıfırlanır. Hedef seyreklik
$s$, eğitim epoch'u $e$ boyunca
\begin{equation}
s_e = s \cdot \min\!\Bigl(1, \tfrac{e}{e_{\text{ramp}}}\Bigr)^{3}
\end{equation}
kübik çizelgesini izler; burada $e_{\text{ramp}} = 50$ olup bunu
50 epoch'luk maske-sabit ince ayar izler.
$s \in \{0.3, 0.5, 0.7, 0.9\}$ değerleri beş tohumla tarandıktan
sonra (Bölüm~\ref{sec:results:sparsity}), doğruluk ile
sıkıştırılabilirliği dengeleyen U-eğrisi optimumu olarak $s = 0.5$
(283 sıfır olmayan parametre) seçilmiştir.

\subsection{Aktivasyon Kalibrasyonlu Tensör-Başına Q15 Nicelendirme}
\label{sec:method:q15}

Her ağırlık tensörü $\mathbf{W}^{(\ell)}$
($\ell \in \{W_1, W_2, U_1, U_2\}$) aşağıdaki biçimde nicelendirilir:
\begin{equation}
\hat{w}^{(\ell)}_{ij}
  = \mathrm{clip}\!\Bigl(
        \mathrm{round}\bigl(w^{(\ell)}_{ij} / s_\ell\bigr),
        -2^{15}, 2^{15}-1
    \Bigr),
\end{equation}
Burada tensöre özgü ölçek $s_\ell$, $\max_{ij}|w_{ij}| / s_\ell$
değerinin Q15 tavanının hemen altında kalacağı biçimde seçilir.
Çıkarım sırasında ileri geçişte çözülmüş değer
$\hat{w}^{(\ell)}_{ij} \cdot s_\ell$ kullanılır.

Bu şemayı aktivasyonlara doğrudan Q15 $[-1, 1)$ biçiminde
uygulamak yıkıcı sonuç verir: eğitilmiş modelimizde gizli durum
$\mathbf{h}_t$ $\sim$62 büyüklüklerine ulaşır; bu değer Q15'in
temsil edilebilir aralığının bir büyüklük mertebesi ötesindedir.
Müdahale edilmediğinde F1 0.918'den 0.16'ya düşer ve
\textsc{standing} sınıfı tamamen kaybolur
(Bölüm~\ref{sec:results:quant}).

Basit bir çözüm, gözlenen $\mathbf{h}_t$ büyüklüklerini kapsayan
$\pm 64$ temsil aralığına sahip sabit Q9.6 sabit nokta biçimine
geçmektir. Ancak bu seçim, her tensörün gerçek dinamik aralığını
göz ardı ederek tüm tensörlerde aynı güvenlik payını kullanır ve
çözünürlükten ödün verir. Bunun yerine \emph{aktivasyon başına
kalibrasyon} kullanıyoruz: kalibrasyon geçişi, eğitim
verisinin beş mini-yığınını FP32 modelden geçirir, her ara tensörün
ampirik maksimumunu kaydeder, \%10 güvenlik payı uygular ve her
aktivasyona kendi ölçeğini atar. Tensör-başına şema Q9.6'yı
genelleştirir; buna ihtiyaç duyan tensörler için (ör.\ $\mathbf{h}_t$)
$\pm 64$ aralığına uyarlanabilir biçimde yaklaşırken, buna ihtiyaç
duymayan tensörlerde (ör.\ kapı çıktısı
$\mathbf{z}_t \in [0,1]$) tam Q15 çözünürlüğünü korur.
Kalibrasyonla birlikte tamsayı C uygulaması ile FP32 referansı
arasındaki tahmin uyumu 3{,}399 pencerelik test kümesinin tamamında
\%100'dür ve C tarafındaki makro F1 0.9176'ya ulaşır; bu değer
makale boyunca raporlanan gömülü model sonucudur.

\subsection{LUT Tabanlı Aktivasyonlar}
\label{sec:method:lut}

Çarpıcısız bir hedefte en pahalı işlemler $\sigma$ ve $\tanh$
hesaplarıdır. Her çağrı yazılımla öykünülen aşkın fonksiyonlara
gider; hücre de örnek başına $2H$ aktivasyon hesaplar ($H{=}16$ için
$2 \cdot 16 = 32$, 128 örnekli bir pencerede 4{,}096 çağrı).
Çalışma zamanındaki bu çağrıları, Flash'ta saklanan ve $[-8, +8]$
giriş aralığında tanımlanan 256 girişli arama tablolarıyla
değiştiriyoruz:
\begin{equation}
\mathrm{LUT}[k]
  = f\!\Bigl( -8 + k \cdot \tfrac{16}{255} \Bigr), \quad
  k \in \{0, 1, \ldots, 255\},
\end{equation}
Burada $f \in \{\sigma, \tanh\}$. $[-8, +8]$ dışındaki girişler
$\pm 1$ değerlerine doyurulur; bu bölgelerde her iki fonksiyon da
kayan nokta duyarlığında zaten bu değerlere çok yakındır. Aralık
içinde ise bitişik girdiler arasında doğrusal ara değerleme yapılır.
Böylece her aktivasyon bir karşılaştırma, iki indeksli yükleme, bir
çıkarma ve bir çarp-topla işlemine iner. İki tablo birlikte 2~KB
Flash kaplar (256 giriş $\times$ 4~bayt $\times$ 2 tablo); bu
değişiklik \msp{} üzerinde uçtan uca
\textbf{30.5$\times$} hızlanma sağlar
(Bölüm~\ref{sec:results:streaming}).

\begin{figure}[t]
\centering
\begin{tikzpicture}[
  node distance=4mm and 4mm,
  every node/.style={font=\scriptsize},
  stage/.style={draw, rounded corners=2pt, align=center,
                 minimum height=8mm, minimum width=14mm, fill=blue!8,
                 inner sep=2pt},
  deploy/.style={draw, rounded corners=2pt, align=center,
                 minimum height=8mm, minimum width=16mm, fill=green!10,
                 inner sep=2pt},
  arr/.style={-{Stealth[length=4pt]}, semithick},
]
\node[stage] (train)  {Kayan nokta\\eğitim};
\node[stage, right=of train]    (lr)    {Düşük rank\\$r_w{=}2$, $r_u{=}8$};
\node[stage, right=of lr]       (sp)    {IHT\\seyreklik 0.5};
\node[stage, right=of sp]       (q15)   {Q15\\+ kalib.};

\node[deploy, below=6mm of lr]  (arduino) {Arduino Uno\\(AVR, $\times$HW)};
\node[deploy, below=6mm of q15] (msp)     {MSP430G2553\\($\times$HW yok)};

\node[align=center, below=2mm of arduino, font=\scriptsize\itshape, text=gray]
   (note1) {taşınabilir C + LUT};

\draw[arr] (train) -- (lr);
\draw[arr] (lr)    -- (sp);
\draw[arr] (sp)    -- (q15);
\draw[arr] (q15.south) -- ++(0,-3mm) -| (arduino.north);
\draw[arr] (q15.south) -- ++(0,-3mm) -| (msp.north);
\end{tikzpicture}
\caption{Uçtan uca \fastgrnn{} sıkıştırma ve gömülü hedefe aktarma boru
hattı. Kayan nokta eğitim $\rightarrow$ düşük rank $\rightarrow$
seyrek $\rightarrow$ kalibre edilmiş Q15 ağırlıklar $\rightarrow$
LUT aktivasyonlu taşınabilir C $\rightarrow$ Arduino ve MSP430
ikilileri.}
\label{fig:pipeline}
\end{figure}

\subsection{Uygulama}
\label{sec:method:impl}

Eğitim PyTorch~2.x ile yapılmıştır. Gömülü C çıkarım motoru
(\texttt{fastgrnn.cpp}), hem \texttt{avr-gcc} (\arduino{})
hem de Code Composer Studio \texttt{msp430-elf-gcc} (\msp{})
altında değiştirilmeden derlenen tek bir $\sim$200 satırlık çeviri
birimidir. Tüm ağırlıklar, tensör-başına ölçekler ve iki aktivasyon
LUT'u ile birlikte Flash içinde \texttt{const int16\_t} dizisi
olarak saklanır. Çalışma zamanı çalışma kümesi (gizli durum
$\mathbf{h}$, kapı $\mathbf{z}$, aday $\tilde{\mathbf{h}}$, çıkış
logitleri ve geçici alan) toplamda $\sim$300~bayttır ve \msp{}'nun
512~B SRAM bütçesine rahatça sığar.
